\documentclass{article}
\usepackage[UTF8]{ctex}
\usepackage{amsmath,amssymb,amsfonts}
\usepackage{enumitem}
\usepackage{geometry}
\geometry{a4paper, margin=1in}
% === 错题本排版优化（不参与题目标识） ===
\setlength{\parindent}{0pt}
\setlength{\parskip}{0.35em plus 0.10em minus 0.08em}
\setlength{\emergencystretch}{3em}
\setlist{itemsep=0.25em, topsep=0.40em, parsep=0pt}
\allowdisplaybreaks[2]
% === 错题本排版优化结束 ===
\begin{document}
\section*{重庆大学2024年数学分析真题}
\section*{62．重庆大学2024年数学分析试卷}

\subsection{题目1}
一．计算数列极限

\[
\lim _{n \rightarrow \infty}\left(\frac{1}{\ln (1+n)-\ln n}-n\right)
\]

\subsection{题目2}
二．计算第一类曲面积分 \(\iint_{S}\left(x^{2}+y^{2}\right) \mathrm{d} S\) ，其中 \(S\) 是立体 \(\sqrt{x^{2}+y^{2}} \leq z \leq 1\) 的边界．

\subsection{题目3}
三．求级数 \(\sum_{n=1}^{\infty} \frac{(-1)^{n}}{n(2 n-1)}\) 的和．

\subsection{题目4}
四．计算第二类曲线积分

\[
\int_{L}\left(y^{2}-z^{2}\right) \mathrm{d} x+\left(2 z^{2}-x^{2}\right) \mathrm{d} y+\left(3 x^{2}-y^{2}\right) \mathrm{d} z .
\]

其中 \(L\) 是平面 \(x+y+z=2\) 与柱面 \(|x|+|y|=1\) 的交线，从 \(z\) 轴正方向往下看去，\(L\) 是顺时针方向。

\subsection{题目5}
五．已知三元函数 \(f(x, y, z)\) 在 \(\mathbb{R}^{3}\) 上具有连续的二阶偏导数，设 \(\mathbb{R}^{3}\) 中光滑简单封闭曲面的全体为\(\sum\) ，对于 \(S \in \sum\) ，用 \(\mathbf{n}\) 表示 \(S\) 的外法线单位方向向量，\(\frac{\partial f}{\partial \mathbf{n}}\) 表示 \(f\) 沿 \(\mathbf{n}\) 的方向导数．\\
（1）证明：\(f\) 在 \(\mathbb{R}^{3}\) 上满足

\[
\frac{\partial^{2} f}{\partial x^{2}}+\frac{\partial^{2} f}{\partial y^{2}}+\frac{\partial^{2} f}{\partial z^{2}}=0 \text { (即 } f \text { 为 } \mathbb{R}^{3} \text { 上的调和函数) }
\]

当且仅当 \(\iint_{S} \frac{\partial f}{\partial n} \mathrm{~d} S=0\left(\forall S \in \sum\right)\) ．\\
（2）设 \(f(x, y, z)\) 为 \(\mathbb{R}^{3}\) 上的调和函数，\(S \in \sum\) ，且 \(S\) 围成的有界闭区域记作 \(V\) ．证明：

\[
f\left(x_{0}, y_{0}, z_{0}\right)=\frac{1}{4 \pi} \iint_{S}\left(f \frac{\cos (\mathbf{r}, \mathbf{n})}{|\mathbf{r}|^{2}}+\frac{1}{|\mathbf{r}|} \frac{\partial f}{\partial \mathbf{n}}\right) \mathrm{d} S
\]

其中 \(\left(x_{0}, y_{0}, z_{0}\right)\) 为 \(V\) 内部的一个定点，\(r=\left(x-x_{0}, y-y_{0}, z-z_{0}\right),(r, n)\) 表示两向量的夹角，\(|r|\) 表示向量 \(\mathbf{r}\) 的模长．\\
（3）若函数 \(f\) 和闭区间［0，1］上的连续函数 \(g\) 在单位球 \(\Omega=\left\{(x, y, z) \mid x^{2}+y^{2}+z^{2} \leq 1\right\}\) 上满足

\[
\frac{\partial^{2} f}{\partial x^{2}}+\frac{\partial^{2} f}{\partial y^{2}}+\frac{\partial^{2} f}{\partial z^{2}}=g\left(x^{2}+y^{2}+z^{2}\right) .
\]

证明： \(\iiint_{\Omega}\left(x \frac{\partial f}{\partial x}+y \frac{\partial f}{\partial y}+z \frac{\partial f}{\partial z}\right) \mathrm{d} x \mathrm{~d} y \mathrm{~d} z=\pi \int_{0}^{1} \sqrt{t}(t-1) g(t) \mathrm{d} t\) ．

\subsection{题目6}
六。判断下列四个命题的正误。正确的，请给出证明；错误的，需举出反例。

\begin{itemize}
  \item[（1）] 若正项级数 \(\sum_{n=1}^{\infty} a_{n}\) 收敛，则 \(\lim _{n \rightarrow \infty} \frac{a_{n+1}}{a_{n}}=1\) 。
\end{itemize}

（2）若函数 \(f(x)\) 在 \(\mathbb{R}\) 上连续且有界，则 \(f(x)\) 在 \(\mathbb{R}\) 上一致连续．\\
（3）若二元函数 \(f(x, y)\) 在 \((0,0)\) 点关于 \(x\) 和 \(y\) 都不存在偏导数，则 \(f(x, y)\) 在点 \((0,0)\) 沿 \(x\) 轴的方向导数必不存在．\\
（4）若函数 \(f(x)\) 在闭区间 \([a, b]\) 上有界且只存在有限个不连续点，则函数 \(f(x)\) 在 \([a, b]\) 上 Riemann可积。

\subsection{题目7}
七．设参数 \(\alpha>0\) ，讨论广义积分 \(\int_{0}^{+\infty} \frac{e^{\sin x} \sin (2 x)}{x^{\alpha}} \mathrm{d} x\) 的收敛性（包括条件收敛和绝对收敛）．

\subsection{题目8}
八．若函数 \(f(x)\) 和 \(g(x)\) 都在闭区间 \([a, b]\) 上连续，在 \((a, b)\) 上二阶可导，并且它们在 \((a, b)\) 内有相等的最大值，\(f(a)=g(a), f(b)=g(b)\) ，求证：存在 \(\xi \in(a, b)\) ，使得 \(f^{\prime \prime}(\xi)=g^{\prime \prime}(\xi)\) 。

\subsection{题目9}
九．函数列 \(\left\{f_{n}(x)\right\}_{n=1}^{\infty}\) 在闭区间 \([a, b]\) 上逐点收敛于 \(f(x)\) ．\\
（1）假设每个 \(f_{n}(x)\) 都在 \([a, b]\) 上连续，并且 \(\left\{f_{n}(x)\right\}_{n=1}^{\infty}\) 一致收敛，证明：\(f(x)\) 在 \([a, b]\) 上连续．\\
（2）假设每个 \(f_{n}(x)\) 都在 \([a, b]\) 上 Riemann 可积，并且 \(\left\{f_{n}(x)\right\}_{n=1}^{\infty}\) 一致收敛，证明：\(f(x)\) 在 \([a, b]\) 上Riemann 可积。\\
（3）假设每个 \(\left\{f_{n}(x)\right\}\) 都在 \([a, b]\) 上可微，存在常数 \(M>0\) ，对任意的 \(x\) 和正整数 \(n\) ，都有 \(\left|f_{n}^{\prime}(x)\right| \leq M\) ．证明：\(\left\{f_{n}(x)\right\}_{n=1}^{\infty}\) 在 \([a, b]\) 上一致收敛。
\end{document}
